% MCFG #2 v74: arbitrary finite-group shortening, Japanese mirror.
\section{固定窓と任意の有限群の組合せ}
\label{app:window-group}

本付録では，\ref{app:window-abelian} の
有限群に関する可換性の仮定を取り除く．
学習作用素および窓のみの場合の上界は変更しない．
有限群とファンアウトは固定パラメータである．
以下の定数は非常に大きいが，重要なのは，
\emph{型付け前の}文法サイズと規則厚みについての多項式性である．

\begin{lemma}[均質なブロック配置の構造]
\label{lem:group-block-order}
$e\ge1,n\ge3$ とする．各 $t\in[e]$ と
$i\in[n]$ に対応する相異なる点 $(t,i)$ が
一直線上に全順序で並んでいるとする．
(i) 同じ添字 $i$ における $(t,i)$ と $(u,i)$ の順序は
$i$ に依存しない．
(ii) 任意の順序付き型対 $(t,u)$（$t=u$ も許す）について，
$i<j$ のときの $(t,i)$ と $(u,j)$ の順序は
$i<j$ の取り方に依存しない．
このとき型の集合は，それぞれ連続した領域を占める
いくつかのクラスに分割される．各クラスには
添字に関する一つの向き（昇順または降順）があり，
領域内はその向きに沿うランク別の連続した
$n$ 個のチャンクに分かれる．
各チャンクはクラス内の各型の点をちょうど一つずつ含み，
同じチャンク内の型の順序はランクによらない．
\end{lemma}
\begin{proof}
(ii) で $t=u$ と置くと，各型の点列は単調である．
相異なる型 $t,u$ に対し，$i<j$ のときの
$(t,i)<(u,j)$ を $X$，$(u,i)<(t,j)$ を $Y$
と書く．(ii) より真偽は添字対に依存しない．
両型が昇順のとき，$X\land Y$ なら
ランク $i$ の両型の点がランク $j>i$ の両型の
点よりすべて左にある．
$X\land\neg Y$，$\neg X\land Y$ では
一方の型の全点が他方の全点より左にある．
$\neg X\land\neg Y$ なら
\[
 (t,1)<(t,2)<(u,1)<(u,2)<(t,1)
\]
という矛盾が生じる．
$t$ が昇順，$u$ が降順のとき，
$X\land Y$ なら
$(t,2)<(u,3)<(u,1)<(t,2)$，
$\neg X\land\neg Y$ なら
$(t,3)<(u,2)<(t,1)<(t,3)$
となって矛盾する．
$X\land\neg Y$ なら異添字の点はすべて
$t$ が左にあり，同添字で逆の比較があると
$(t,1)<(u,2)<(u,1)<(t,1)$
と矛盾する（(i) を用いた）．したがって
全 $t$-点は全 $u$-点の左にある．
$\neg X\land Y$ も対称的に処理できる．
両型が降順の場合は添字順を反転して昇順の場合に帰着する．
ゆえに二型は，全点が分離しているか，
同じ向きでランクごとに交互配置されている．

交互配置関係は推移的である．
$t$ と $u$，$u$ と $v$ が交互配置なら
三型は同方向である．添字を必要に応じて反転し
昇順とする．$t$ と $v$ が分離し，
$t$ が左なら
\[
 (v,1)<(u,2)<(t,3)<(v,1)
\]
と矛盾する．逆に $v$ が左なら
$(t,1)<(u,2)<(v,3)<(t,1)$
と矛盾する．
したがって交互配置関係に等号を加えると同値関係となる．
あるクラスの外の型は，クラス内のすべての型と分離しており，
それらのすべてに対し同じ側に位置する．
異なる側に位置すれば，クラス内の二つの交互配置型が
互いに分離してしまうからである．
よってクラス領域は連続する．
ランク間の順序関係がチャンクの連続性を与え，
(i) が各チャンク内の順序を与える．
\end{proof}

\begin{lemma}[有限群における複数成分同時の区間短縮]
\label{lem:group-interval}
$A$ を任意の有限群とする．$n$ 個の添字が，
高々 $f$ 個の順序付き出力成分へ，
添字一つにつき高々 $f$ 個の重ならないブロックを
寄与するとする．各ブロックには代替ブロックが固定され，
添字を選択するとその添字のすべてのブロックを同時に置換する．
その他の材料は変更しない．
ブロックの出力先と同じ添字内の順序が一定であり，
さらに $i<j$ の型対間の左右比較も添字対によらず一定とする．
$n\ge\max\{3,|A|^f\}$ なら，非空の連続添字区間が存在し，
その添字のブロックを同時置換しても，
すべての出力成分の $A$-値を保存できる．
\end{lemma}
\begin{proof}
補題~\ref{lem:group-block-order} を各根成分に適用する．
全根成分を通じたクラス領域数は $q\le f$ である．
一つのクラス領域 $K$ を考え，チャンクを左から
$1,\ldots,n$ と番号付けする．
ブロック間にある不変の材料も領域の値に含める．
チャンク $p$ の末尾までの元の群積を $P_p$，
チャンク $1,\ldots,p$ だけを置換した場合の
同じ接頭部分の群積を $Q_p$ とし，
$P_0=Q_0=1_A$ と置く．
\[
 \kappa^K_p=P_pQ_p^{-1}
\]
と定義する．
チャンク $p+1,\ldots,p'$ のみを置換した領域全体の積は，
積の順序を維持すると
\[
 P_p\,Q_p^{-1}Q_{p'}\,P_{p'}^{-1}P_n
\]
となる．
実際，$Q_p^{-1}Q_{p'}$ が置換後の該当区間の積であり，
$P_{p'}^{-1}P_n$ が変更されない後続部分の積である．
したがって
$\kappa^K_p=\kappa^K_{p'}$ なら
領域全体の値は元の $P_n$ に等しい．
ここでは逆元しか使わず，積の交換は行っていない．

クラス $K$ の方向が昇順なら添字区間 $(a,b]$ は
チャンク区間 $(a,b]$ に対応し，
降順なら $(n-b,n-a]$ に対応する．
$j=0,\ldots,n$ について，昇順クラスから
$\kappa^K_j$，降順クラスから
$\kappa^K_{n-j}$ を集めて $A^q$ の状態を作る．
$n+1>|A|^q$ なので，
状態が一致する $a<b$ が存在する．
添字 $a+1,\ldots,b$ を同時置換すると
すべてのクラス領域の群積が保存される．
領域外の材料は不変なので，
すべての根成分の群値が保存される．
\end{proof}

\begin{lemma}[重い部分木の反鎖の一様な短縮]
\label{lem:group-heavy-antichain}
\[
 N_{A,f}=\max\{3,|A|^f\},\quad
 c_f=2^{f^2},\quad U_f=f^{f+1}f!,\quad
 M_{A,f}=U_f\,c_f^{\,c_f(N_{A,f}-1)+1}
\]
と置く．$G$ をreduced good MCFG，
$g:\Sigma^*\to A$ を任意の有限群への準同型とする．
根の成分別 $g$-値を指定して，
総出力長，次いで節点数が辞書式最小となる終端導出において，
総出力長が $\bar\theta_G$ より大きい
部分木の反鎖の大きさは $M_{A,f}$ 未満である．
\end{lemma}
\begin{proof}
線形・非削除規則により，部分木の各成分は
根の一つの成分の連続ブロックとして一度だけ出現する．
比較不能な部分木どうしのブロックは重ならない．
反鎖の各節点に対して，同じ\emph{型付け前の}
非終端記号からの長さ高々 $\bar\theta_G$ の
代替終端導出を固定できる
（補題~\ref{lem:thickness-local-bounds}(ii)）．
群の値は変化してもよい．

一節点のブロック数は高々 $f$ であり，
その単項パターン（ブロック数，根成分の行き先，
同一根成分内での自ブロック順序）は高々 $U_f$ 通りである．
まず単項パターンを一致させる．
そのような節点 $i<j$ について，
同じ根成分へ行くすべての型対 $(t,u)$
（$t=u$ を含む）の左右比較を色とする．
色数は高々 $c_f$ である．
初等的な多色Ramsey上界
\[
 R_{c_f}(N)\le c_f^{\,c_f(N-1)+1}
\]
を使う．$c=c_f$，$L=c(N-1)+1$ と置く．
初めに少なくとも $c^L$ 頂点を用意し，
$v_1,\ldots,v_L$ を順に選び，各段階で
後続頂点の最大単色クラスに候補を絞る．
$j<L$ 段階後の候補集合 $S_j$ について
$|S_j|\ge\lceil(|S_{j-1}|-1)/c\rceil
\ge c^{L-j}$ となる．
最初の $L-1=c(N-1)$ 頂点のうち，
同じ出色を持つものが少なくとも $N-1$ 個ある．
最後の $v_L$ は先の各候補集合に属するので，
この色の $N$ 頂点の単色完全部分グラフが得られる．
したがって，大きさ $M_{A,f}$ 以上の反鎖には，
単項パターンと型対の比較色が一定の
$N=N_{A,f}$ 個の節点が存在する．

それらのブロック配置には
補題~\ref{lem:group-block-order} が各根成分で適用できる．
続いて補題~\ref{lem:group-interval} により，
非空の節点区間を同時に短い導出へ置き換えても
根の\emph{すべての群値}は保存される．
もとの節点は重く，代替の各節点の総長は
$\bar\theta_G$ 以下なので，根の総出力長は真に減る．
これは選択した導出の最小性に反する．
\end{proof}

\begin{theorem}[任意の有限群値を指定した短いタプル]
\label{thm:group-typed-shortest}
$G=(V,\Sigma,P,S,\fan)$ をfan-out $\le f$，
規則ランク $\le r_{\max}$（$r_{\max}\ge1$）の
reduced good MCFGとする．
$g:\Sigma^*\to A$ を\emph{任意の有限群}
$A$ への準同型とする．各 $B\in V$ と
実現可能な $\mathbf s\in A^{\fan(B)}$ に対し，
\[
 \tuple u\in L_B(G),\qquad
 g^{(\fan(B))}(\tuple u)=\mathbf s
\]
を満たし，かつ
\[
 \|\tuple u\|_1\le
 \beta_g^{\mathrm{grp}}:=
 (r_{\max}+1)\bar\theta_G
 \max\{1,(M_{A,f}-1)|V||A|^f\}
\]
となるタプルが存在する．とくに群と
$f,r_{\max}$ を固定すると
$\beta_g^{\mathrm{grp}}
=O_{A,f,r_{\max}}(\|G\|\bar\theta_G)$ である．
\end{theorem}
\begin{proof}
$B$ から指定した群ベクトルを実現する導出のうち，
（総長，節点数）が辞書式最小のものを選ぶ．
各節点を非終端記号と各成分の群値で標識する．
祖先と真の子孫で同じ標識を持つことはできない．
祖先の部分木を子孫の部分木に取り換えれば，
規則テンプレートを順序どおり評価することで
根までの群値は変わらず，非削除性から総長は増えず，
節点数は真に減るからである．
よってすべての根から葉までの道の節点数は
\[
 n_g=|V||A|^f
\]
以下である．

総出力長が $\bar\theta_G$ を超える節点を重いと呼ぶ．
重い節点は祖先について上に閉じている．
補題~\ref{lem:group-heavy-antichain} より
重い葉は高々 $M_{A,f}-1$ 個であり，
各重い節点はそのいずれかへの道上にある．
したがって重い節点の総数は高々
$(M_{A,f}-1)n_g$ である．
根が軽ければ長さはすでに $\bar\theta_G$ 以下である．
そうでなければ，各重い節点の規則が明示的に出力する
終端数は高々 $\bar\theta_G$，
その高々 $r_{\max}$ 個の軽い子部分木の総長も
各々高々 $\bar\theta_G$ である．
重い節点と極大軽部分木を数えると結論が従う．
\end{proof}

\begin{lemma}[窓と任意有限群の型付けに対する短い正準タプル]
\label{lem:window-group-short}
定理~\ref{thm:group-typed-shortest} の仮定のもと，
$k,\ell\ge0$ と
$h\preceq(h_\Sigma^{k,\ell},g)$ を満たす
固定有限型付け $h$ を選ぶ．
$G$ の刈り込み済み完全 $h$-精密化を
$\widetilde G_0$ とし，その開始記号以外の
非終端記号数を $n_{\mathrm{typ}}$ とする．
各非開始記号 $X$ について，
\[
 \|\omega(X)\|_1\le
 f(k+\ell+1)n_{\mathrm{typ}}
 \bigl(\bar\theta_G+r_{\max}\beta_g^{\mathrm{grp}}\bigr)
\]
を満たす正準タプルを選べる．
\end{lemma}
\begin{proof}
$X$ からの終端導出で節点数が最小のものを選ぶ．
根から葉への道上に同じ型付き非終端記号が
二度現れれば，根の $h$-型を保存したまま
節点数を減らせるため，各道は
$n_{\mathrm{typ}}$ 節点以下である．
補題~\ref{lem:boundary-short-canonical} と同様に，
各根成分で $h_\Sigma^{k,\ell}$ の値を決める
位置に印を付ける．短い成分では全位置，
長い成分では先頭 $k$ 位置と末尾 $\ell$ 位置を，
$k=\ell=0$ なら一点を保護する．
高々 $f(k+\ell+1)$ 個の印への祖先経路が作る
骨格の節点数は高々
$f(k+\ell+1)n_{\mathrm{typ}}$ である．

骨格から外れる各子部分木の，
元の非終端記号 $B$ と実際の成分別群値を記録する．
定理~\ref{thm:group-typed-shortest} により，
同じ $B$ から，同じ各成分群値を持ち，
総長が $\beta_g^{\mathrm{grp}}$ 以下の代替導出を選べる．
$G$ は良形なので代替成分はすべて非空である．
保護位置は変化せず，長い根成分は長いままであるため，
窓の値は保存される．また，各子成分の群値を
個別に保存しているので，規則テンプレートを
元の順序で評価すれば根の各群値も保存される．
可換性は必要ない．
したがって積型付け $(h_\Sigma^{k,\ell},g)$，
およびより粗い根の $h$-型は不変である．
命題~\ref{prop:output-invariants}(ii) により，
元の $G$-導出を，根を $X$ とする
刈り込み済み完全精密化の導出へ持ち上げられる．

各骨格節点が明示する終端は高々
$\bar\theta_G$ 個であり，
高々 $r_{\max}$ 個の省略子それぞれは
長さ高々 $\beta_g^{\mathrm{grp}}$ を寄与する．
ゆえに上界が従う．
\end{proof}

\begin{theorem}[窓と任意有限群による型付けの多項式特徴データ]
\label{thm:window-group-polydata}
$f,r_{\max},k,\ell,\Sigma$，任意の有限群 $A$，
準同型 $g:\Sigma^*\to A$ および
$h\preceq(h_\Sigma^{k,\ell},g)$ を満たす固定有限型付け
$h$ を固定する．
$\Cmcf{f}{h}$ の言語を生成する任意の
reduced good fan-out-$\le f$，
rank-$\le r_{\max}$ 文法 $G$ に対して，
ランク制限つき再構成作用素
$\mathcal B^{\le r_{\max}}_{f,h}$ は
次を満たす特徴標本を持つ：
\[
 |\Wit(\widetilde G_0)|=O(\|G\|),\qquad
 \|\Wit(\widetilde G_0)\|_+
 =O_{f,r_{\max},k,\ell,h,g,\Sigma}
   (\|G\|^3\bar\theta_G).
\]
\end{theorem}
\begin{proof}
補題~\ref{lem:window-group-short} より，
刈り込み済み完全精密化の正準タプルについて
\[
 b_{\mathrm{can}}\le
 f(k+\ell+1)n_{\mathrm{typ}}
 (\bar\theta_G+r_{\max}\beta_g^{\mathrm{grp}})
 =O(\|G\|^2\bar\theta_G)
\]
を得る．ここで型付けと規則ランクの上界は固定され，
$n_{\mathrm{typ}}\le|V||M|^f=O(\|G\|)$ である．
補題~\ref{lem:untyped-contexts} より
各特徴証拠語の長さは
$2\bar\theta_G+r_{\max}b_{\mathrm{can}}
=O(\|G\|^2\bar\theta_G)$ 以下となり，
命題~\ref{prop:witness-cardinality} より
語数は $O(\|G\|)$ 以下である．
空語の証拠を加えても総符号化長は定数しか増えない．
\end{proof}

\begin{remark}
この強化は，任意の有限\emph{モノイド}型付けについて
元文法の厚さから多項式上界が得られると主張しない．
定数 $M_{A,f}$ は粗いRamsey上界である．
$A$ が可換な場合は，\ref{app:window-abelian} の
直接証明の方が小さい定数を与える．
\end{remark}
